\section{Initial separation and the geometric criterion}
\label{sec:geometry}

Let $\Phi_h$ denote a flow map for $h\geq0$, and let
$B^d=\{\xi\in\R^d:\norm{\xi}_2\leq1\}$. Write $\sigma$ for the
effective width $\sigma_F$ in the source's geometric definition.
Work in a coordinate chart in which $B(x,\sigma)$ and the short trajectories
considered below remain inside the domain, so coordinate wrapping is absent.
The maximal relative drift in source Eq.~(30), p.~5, is
\begin{equation}
 D_i^{\rm raw}(h;x,\sigma)
 :=\max_{\xi\in B^d}
 \left|\Phi_h(x+\sigma\xi)_i-\Phi_h(x)_i\right|.
 \label{eq:raw-separation}
\end{equation}
It is a separation of final positions. Source Eq.~(30) then identifies it with
$h\max_{\xi\in B^d}|F_i(x+\sigma\xi)-F_i(x)|+O(h^2)$.
That identification drops the separation already present at time zero.
Indeed, a uniform short-time flow expansion gives
\begin{equation}
 \Phi_h(x+\sigma\xi)_i-\Phi_h(x)_i
 =\sigma\xi_i+h\bigl(F_i(x+\sigma\xi)-F_i(x)\bigr)+O(h^2).
 \label{eq:flow-offset}
\end{equation}
For fixed $\sigma>0$, the first term cannot be included in an $O(h^2)$
remainder. In particular, $D_i^{\rm raw}(0;x,\sigma)=\sigma$.

\begin{proposition}[Obstruction for the stated directional maps]
\label{prop:raw-obstruction}
Suppose $\partial_iF_i=0$ on the relevant neighborhood. The directional
flow used in a Trotter substep is exactly
\begin{equation}
 \Phi_h^{(i)}(x)=x+hF_i(x)e_i.
 \label{eq:directional-map}
\end{equation}
If $D_i^{\rm raw}$ is evaluated using this map, then
$D_i^{\rm raw}\geq\sigma$. Consequently, if each component in source
Eq.~(31) is evaluated using its corresponding directional substep,
\begin{equation}
 \frac{\bigl(\sum_{i=1}^d(D_i^{\rm raw})^2\bigr)^{1/2}}
 {\sqrt d\,\sigma}\geq1.
 \label{eq:ratio-obstruction}
\end{equation}
\end{proposition}

\begin{proof}
Choose $\xi=e_i$ in \eqref{eq:raw-separation}. The coefficient does not
change under this perturbation because it is independent of coordinate
$i$. Hence
$\Phi_h^{(i)}(x+\sigma e_i)_i-\Phi_h^{(i)}(x)_i=\sigma$.
Taking the maximum proves the componentwise bound and summing proves
\eqref{eq:ratio-obstruction}. Centers at different intermediate substeps
do not change the argument, provided their neighborhoods remain in the chart.
\end{proof}

Thus the approximately $0.1$ condition printed in source Eq.~(31) cannot
hold under the directional-map interpretation given in its geometric
discussion. If the intended map were instead the complete flow, the
time-zero contradiction in \eqref{eq:flow-offset} would still remain:
the normalized raw separation tends to $1$, rather than $0$, as $h\to0$.
The following nonlinear example also gives an all-$h$ obstruction for the
complete-flow interpretation, rather than only a time-zero objection.

This obstruction occurs for a genuinely nonlinear conservative field. Set
\begin{equation}
 F(q,p)=(\sin(2\pi p),0),\qquad
 H(q,p)=-\frac{\cos(2\pi p)}{2\pi}.
 \label{eq:nonlinear-shear}
\end{equation}
Then $\dot q=\partial_pH$, $\dot p=-\partial_qH$, and
$\Phi_h(q,p)=(q+h\sin(2\pi p),p)$. Both own-coordinate derivatives
vanish. Here the complete flow equals the sole nonidentity directional
factor exactly; the other factor is the identity. Thus
\eqref{eq:ratio-obstruction} applies to both interpretations for every
admissible $h$. A packet with means $(1/2,1/2)$ and a $p$ marginal symmetric
about $1/2$ has $\E{\sin(2\pi P)}=0$. Its mean therefore remains at the
fixed-point trajectory $(1/2,1/2)$, exactly, on any interval without
coordinate wrapping. A smooth compactly supported symmetric packet can
ensure this condition. The printed geometric condition nevertheless fails.
This example concerns the geometric condition; its covariance need not
remain constant, and no constant-covariance error theorem is invoked.

\subsection{A corrected displacement quantity}
An explicit repair is to compare the increments, subtracting the initial
separation inside the absolute value:
\begin{equation}
 D_i^{\rm inc}(h;x,\sigma):=
 \max_{\xi\in B^d}\left|
 [\Phi_h(x+\sigma\xi)-(x+\sigma\xi)]_i
 -[\Phi_h(x)-x]_i\right|.
 \label{eq:increment-distortion}
\end{equation}
For the exact directional map \eqref{eq:directional-map},
\begin{equation}
 D_i^{\rm inc}
 =h\max_{\xi\in B^d}|F_i(x+\sigma\xi)-F_i(x)|.
 \label{eq:increment-exact}
\end{equation}
This equality has no time-discretization remainder. It uses the
own-coordinate condition and is specific to a directional substep.

\begin{proposition}[Local increment bound]
\label{prop:increment-bound}
Assume $F_i\in C^2$ on $B(x,\sigma)$ and use the directional map
\eqref{eq:directional-map}. Then
\begin{equation}
 \frac{D_i^{\rm inc}}{\sigma}
 \leq h\norm{\nabla F_i(x)}_2+
 \frac{h\sigma}{2}
 \sup_{y\in B(x,\sigma)}\norm{D^2F_i(y)}_{\rm op}.
 \label{eq:increment-bound}
\end{equation}
If $F_i\in C^3$ and its third derivative has operator norm at most $M_3$
on the ball, the supremum-Hessian term may instead be replaced by
\begin{equation}
 \frac{h\sigma}{2}\norm{D^2F_i(x)}_{\rm op}
 +\frac{h\sigma^2M_3}{6}.
 \label{eq:increment-third-order}
\end{equation}
\end{proposition}

\begin{proof}
Apply the one-dimensional Taylor formula along the segment
$s\mapsto x+s\sigma\xi$. The linear term has absolute value at most
$\sigma\norm{\nabla F_i(x)}_2$ and the integral quadratic remainder is
at most $\sigma^2\sup\norm{D^2F_i}_{\rm op}/2$.
Substitute in \eqref{eq:increment-exact}. Expanding one order further gives
\eqref{eq:increment-third-order}.
\end{proof}

\subsection{A width-uniform complete-flow remainder}
The remainder for a complete flow can be bounded uniformly as the packet
width shrinks. A bound on two separate trajectory remainders would give a
loose $O(h^2/\sigma)$ estimate after normalization. Keeping their difference
inside the integral instead gives the following standard auxiliary ODE
estimate. It concerns the exact complete flow, rather than a global error
theorem for a Trotter approximation.

\begin{proposition}[Complete-flow increment bound]
\label{prop:complete-increment}
Let $U\subset\R^d$ be open and convex, and let $F\in C^2(U;\R^d)$.
Suppose $B(x,\sigma)\subset U$, $\sigma>0$, and all trajectories starting
in that ball remain in $U$ for $0\leq t\leq h_0$. Assume the uniform bounds
\begin{equation}
 \sup_U\norm{F}_2\leq M,\qquad
 \sup_U\norm{DF}_{\rm op}\leq J,\qquad
 \sup_U\norm{D^2F}_{\rm bil}\leq H_{\rm v},\qquad
 K:=H_{\rm v}M+J^2,
 \label{eq:complete-constants}
\end{equation}
where $\norm{D^2F}_{\rm bil}$ is the operator norm of the vector-valued
bilinear map $(u,v)\mapsto D^2F[u,v]$. Put
$A_h(z)=\Phi_h(z)-z$ and
\begin{equation}
 R_h(x,\delta):=A_h(x+\delta)-A_h(x)
              -h\bigl(F(x+\delta)-F(x)\bigr).
 \label{eq:complete-remainder-definition}
\end{equation}
Then, for $0\leq h\leq h_0$ and $\norm{\delta}_2\leq\sigma$,
\begin{equation}
 \norm{R_h(x,\delta)}_2
 \leq \frac{K h^2e^{Jh}}{2}\norm{\delta}_2.
 \label{eq:complete-remainder-bound}
\end{equation}
Consequently, writing
$H_i=\sup_{B(x,\sigma)}\norm{D^2F_i}_{\rm op}$, the complete-flow
increment statistic satisfies
\begin{equation}
 \frac{D_i^{\rm inc}(h;x,\sigma)}{\sigma}
 \leq h\norm{\nabla F_i(x)}_2
       +\frac{h\sigma H_i}{2}
       +\frac{K h^2e^{Jh}}{2}.
 \label{eq:complete-increment-bound}
\end{equation}
\end{proposition}

\begin{proof}
Set $G=DF\,F$. Differentiating $G$ gives
$DG[v]=D^2F[v,F]+DF\,DF\,v$, so
$\sup_U\norm{DG}_{\rm op}\leq K$. Convexity permits the corresponding
Lipschitz estimate for $G$. Twice integrating the ODE yields
\begin{equation}
 A_h(z)=hF(z)+\int_0^h\!\int_0^s G(\Phi_r(z))\,\dd r\,\dd s.
 \label{eq:complete-double-integral}
\end{equation}
Subtract the formulas for $z=x+\delta$ and $z=x$ before taking norms.
The flow Lipschitz estimate
$\norm{\Phi_r(x+\delta)-\Phi_r(x)}_2\leq e^{Jr}\norm{\delta}_2$
therefore gives
\[
 \norm{R_h(x,\delta)}_2
 \leq K\norm{\delta}_2\int_0^h\!\int_0^s e^{Jr}\,\dd r\,\dd s
 \leq\frac{K h^2 e^{Jh}}2\norm{\delta}_2.
\]
Combining this with the segmentwise second-order Taylor bound for
$F_i(x+\delta)-F_i(x)$ proves \eqref{eq:complete-increment-bound}.
\end{proof}

The constants in \eqref{eq:complete-remainder-bound} are independent of
$\sigma$ when the same neighborhood and derivative bounds apply. In a
coupled regime $\sigma=\kappa h$ with fixed $\kappa>0$,
\begin{equation}
 \frac{D_i^{\rm inc}}{\sigma}
 \leq h\norm{\nabla F_i(x)}_2
 +\frac{\kappa H_i h^2}{2}+\frac{K h^2e^{Jh}}2.
 \label{eq:coupled-width-step}
\end{equation}
The nonlinear-curvature and complete-flow temporal terms are both
$O(h^2)$ and must both be retained. This local estimate does not establish
a global trajectory-reconstruction criterion.

An additional numerical group uses the nontrivial complete Hamiltonian
flow $F(q,p)=(\sin(2\pi p),\sin(2\pi q))$, with
$H=-\cos(2\pi p)/(2\pi)+\cos(2\pi q)/(2\pi)$. On $\R^2$, valid
global constants are $M=\sqrt2$, $J=2\pi$, and
$H_{\rm v}=(2\pi)^2$. The script compares the normalized remainder with
\eqref{eq:complete-remainder-bound} for $h\in\{.01,.02,.04\}$ and
$\sigma=\kappa h$, $\kappa\in\{.5,1,2\}$. It integrates the center
and 128 boundary perturbations by RK4, comparing 1024 with 2048 internal
steps. These are converged numerical checks of selected trajectories,
not exact continuum executions or a certification of all perturbations;
the uniform statement follows from the proposition's proof.

The two terms in \eqref{eq:increment-bound} separate linear stretching
from local curvature. Increment distortion itself can be nonzero for an
affine field, even though averaging then follows the exact trajectory.
For a directional substep, if only the nonlinear part is desired, one can further subtract
$h\sigma\nabla F_i(x)\cdot\xi$ inside \eqref{eq:increment-distortion}; its
normalized magnitude is bounded by the Hessian term alone. For a complete
flow, the additional temporal remainder in
\eqref{eq:complete-increment-bound} must also be included.
These are standard local Taylor estimates. They repair the definition and
expose its assumptions; they do not establish a new universal
Ehrenfest--Reynolds threshold or a global trajectory-accuracy theorem.
